% Zone „Das kennst du schon“ – aus bank/brueche-dezimalzahlen/zone.jsonl (zusammenbau v0.1)
\einheitenkopf[zone][Kennst du schon]{Das kennst du schon}
\setcounter{aufgabe}{0}

%% TODO Titel: Ich-kann-Satz fehlt in der Bank (Platzhalter: Kettenname) – Nr. 1
%% TODO Anweisung: ein Satz über der ersten Teilaufgabe fehlt in der Bank – Nr. 1
\begin{aufgabe}{Teilen und Vervielfachen im Kopf (vierundzwanzig geteilt durch sechs, drei mal neun), Einmaleins}
\begin{teile}
\teil $56 : 7$
\teil $4{,}2 : 6$
\end{teile}
\end{aufgabe}

%% TODO Titel: Ich-kann-Satz fehlt in der Bank (Platzhalter: Kettenname) – Nr. 2
%% TODO Anweisung: ein Satz über der ersten Teilaufgabe fehlt in der Bank – Nr. 2
\begin{aufgabe}{Größen mit Komma lesen und umrechnen (1,2 kg = 1200 g; 2,50 €)}
\begin{teile}
\teil $3{,}5$ kg – wie viel Gramm? \leerfeld[g]
\teil $0{,}85$ l – wie viel Milliliter? \leerfeld[ml]
\end{teile}
\end{aufgabe}

%% TODO Titel: Ich-kann-Satz fehlt in der Bank (Platzhalter: Kettenname) – Nr. 3
%% TODO Anweisung: ein Satz über der ersten Teilaufgabe fehlt in der Bank – Nr. 3
\begin{aufgabe}{Kreissektor als Anteil vom Vollkreis (120° von 360°)}
\begin{teile}
\teil $180^\circ$ von $360^\circ$ – welcher Anteil vom Vollkreis?
\teil $45^\circ$ von $360^\circ$ – welcher Anteil vom Vollkreis?
\end{teile}
\end{aufgabe}

%% TODO Titel: Ich-kann-Satz fehlt in der Bank (Platzhalter: Kettenname) – Nr. 4
%% TODO Anweisung: ein Satz über der ersten Teilaufgabe fehlt in der Bank – Nr. 4
\begin{aufgabe}{Teiler und Vielfache einer Zahl angeben (Teiler von zwölf; Vielfache von vier)}
\begin{teile}
\teil Nenne alle Teiler von $10$.
\teil Nenne alle Teiler von $36$.
\end{teile}
\end{aufgabe}

%% TODO Titel: Ich-kann-Satz fehlt in der Bank (Platzhalter: Kettenname) – Nr. 5
%% TODO Anweisung: ein Satz über der ersten Teilaufgabe fehlt in der Bank – Nr. 5
\begin{aufgabe}{Stellenwerte natürlicher Zahlen und Zahlenstrahl mit natürlichen Zahlen (Skala ablesen)}
\begin{teile}
\teil Welchen Stellenwert hat die $7$ in $4\,782$?
\teil Von $0$ bis $1\,000$ sind $10$ gleich lange Abschnitte. Welche Zahl gehört an den dritten Teilstrich nach der $0$?
\end{teile}
\end{aufgabe}

%% TODO Titel: Ich-kann-Satz fehlt in der Bank (Platzhalter: Kettenname) – Nr. 6
%% TODO Anweisung: ein Satz über der ersten Teilaufgabe fehlt in der Bank – Nr. 6
\begin{aufgabe}{Schriftlich oder mit Taschenrechner dividieren (drei geteilt durch acht)}
\begin{teile}
\teil $84 : 4$
\teil $7 : 4$
\end{teile}
\end{aufgabe}

%% TODO Titel: Ich-kann-Satz fehlt in der Bank (Platzhalter: Kettenname) – Nr. 7
%% TODO Anweisung: ein Satz über der ersten Teilaufgabe fehlt in der Bank – Nr. 7
\begin{aufgabe}{Prozent als Hundertstel (45 \% = 0,45)}
\begin{teile}
\teil $30\,\%$ – als Bruch mit dem Nenner $100$?
\teil $1{,}2$ – wie viel Prozent? \leerfeld[\%]
\end{teile}
\end{aufgabe}

%% TODO Titel: Ich-kann-Satz fehlt in der Bank (Platzhalter: Kettenname) – Nr. 8
%% TODO Anweisung: ein Satz über der ersten Teilaufgabe fehlt in der Bank – Nr. 8
\begin{aufgabe}{Prozent als Hundertstel (45 \% = 0,45) – Fehler finden}
\begin{teile}
\teil Lea schreibt: $6\,\% = 0{,}6$. Finde den Fehler.
\end{teile}
\end{aufgabe}

%% TODO Titel: Ich-kann-Satz fehlt in der Bank (Platzhalter: Kettenname) – Nr. 9
%% TODO Anweisung: ein Satz über der ersten Teilaufgabe fehlt in der Bank – Nr. 9
\begin{aufgabe}{Prozent als Hundertstel (45 \% = 0,45) – selbst rechnen}
\begin{teile}
\teil $9\,\%$ – als Kommazahl?
\end{teile}
\end{aufgabe}
